\Needspace{8\baselineskip}\originalchapter{2002 年期中卷：冷波、漂移与箍缩}
2002 Fall 22.611J 期中共原题1—6，各有配对扫描官方解答。下文按原题为六个综合组，不把简答小问拆成新题；正文推导为编者独立复核。
\begin{exercise}\label{ex:hist02-mid-1}
原题1（21分）：写电子等离子体频率、离子回旋频率、Debye 长度及 MHD 静力平衡式；判断：(d) Debye 屏蔽使远小于 $\lambda_D$ 的尺度准中性；(e) 电子等离子振荡仅在远小于 $\lambda_D$ 的尺度才有电场和非中性；(f) 良定义的等离子体 $\Lambda=n\lambda_{De}^3$ 很大。
\end{exercise}
\begin{solution}
编者独立解答。$\omega_{pe}^2=n_ee^2/(\epsilon_0m_e)$，$\Omega_i=Z_ieB/m_i$（若要求方向则保留电荷符号），$\lambda_{De}^2=\epsilon_0T_e/(n_ee^2)$，平衡为 $\nabla p=\mathbf J\times\mathbf B$。判断依序：d 假，屏蔽只在尺度大于 $\lambda_D$ 时压低电荷分离；e 假，均匀 Langmuir 振荡可在 $k\lambda_D\ll1$ 发生，只是相对密度扰动小而电场非零；f 真，$\Lambda\gg1$ 表示 Debye 体积粒子多、弱离散涨落。$\lambda_D$ 公式以温度能量单位计。
\end{solution}

\begin{exercise}\label{ex:hist02-mid-2}
原题2（16分）：横波色散函数 $D=1-\omega_{pe}^2/\omega^2-k^2c^2/\omega^2$。画正频 $\omega(k)$，给波能、截止频率，解释低于截止的外驱波。
\end{exercise}
\begin{solution}
编者独立解答。$D=0$ 给 $\omega_+(k)=\sqrt{\omega_{pe}^2+c^2k^2}$，偶于 $k$，在 $k=0$ 起点 $\omega_{pe}$，大 $k$ 渐近 $c|k|$。冷等离子体 $\epsilon=1-\omega_{pe}^2/\omega^2$，电磁波时均波能
\begin{equation}W=\frac{\epsilon_0|E|^2}{4}\frac{\partial(\omega\epsilon)}{\partial\omega}+\frac{|B|^2}{4\mu_0}=\frac{\epsilon_0|E|^2}{2}\quad\text{在色散根上},\end{equation}
其中第一项含电子机械振动能。固定实 $\omega<\omega_{pe}$ 时 $k=\ii\sqrt{\omega_{pe}^2-\omega^2}/c$，半无限介质只有指数衰减的倏逝场，不存在向内传播的体波；有限厚度仍可能有耦合穿透。
\end{solution}

\begin{exercise}\label{ex:hist02-mid-3}
原题3（16分）：两带电粒子 Rutherford 散射在大小碰撞参数截断后给“总 Coulomb 截面”；估算 $T_e$ 电子对离子的动量碰撞频率，写稳态电子动量式以求电场驱动速度及电导率。
\end{exercise}
\begin{solution}
编者独立解答。设相对速度 $v$、约化质量 $\mu$、$b_{90}=|q_1q_2|/(4\pi\epsilon_0\mu v^2)$，有 $\tan(\theta/2)=b_{90}/b$。计数用的总几何截面 $\sigma_{\rm count}=\pi(b_{\max}^2-b_{\min}^2)$ 强依远截断；输运所需是
\begin{equation}\sigma_{\rm mt}=\int_{b_{\min}}^{b_{\max}}(1-\cos\theta)2\pi b\dd b=2\pi b_{90}^2\ln\frac{b_{\max}^2+b_{90}^2}{b_{\min}^2+b_{90}^2}\simeq4\pi b_{90}^2\ln\Lambda.
\end{equation}
$\nu_{ei}\sim n_iv\sigma_{\rm mt}\propto n_iZ_i^2e^4\ln\Lambda/(\epsilon_0^2m_e^{1/2}T_e^{3/2})$。稳态 $0=-en_e\mathbf E-m_en_e\nu_{ei}\mathbf u_e$（离子静止、忽略压强梯度），$\mathbf J=-en_e\mathbf u_e=\sigma\mathbf E$，$\sigma=n_ee^2/(m_e\nu_{ei})$。速度平均需完整分布和碰撞算子才能定精确 Spitzer 系数；不能把 $\sigma_{\rm count}$ 直接代入动量频率。
\end{solution}

\begin{exercise}\label{ex:hist02-mid-4}
原题4（16分）：纯环向磁场的环形容器（原图示 $\mathbf B$ 出纸面），分析离子、电子在外侧磁梯度与曲率下的漂移、由此电荷分离的电场及其 $E\times B$ 共同流，说明为何不能闭合约束；附加题要求定量估计及极化。
\end{exercise}
\begin{solution}
编者独立解答。$B_\phi\propto1/R$，$\mathbf v_{\nabla B}=\mu\mathbf B\times\nabla B/(qB^2)$，$\mathbf v_c=mv_\parallel^2\mathbf B\times\boldsymbol\kappa/(qB^2)$。两项对正负电荷方向相反、量级分别 $T_\perp/(|q|BR)$ 与 $2T_\parallel/(|q|BR)$（能量记法），故上下电荷分离，生竖向 $\mathbf E$；随后两种电荷的 $\mathbf v_E=\mathbf E\times\mathbf B/B^2$ 同向横越磁面，沿环形截面由内向外或反向随图的 $\mathbf B$ 方向而定。物理关键是它有径向非零分量且没有极向场线转动来平均两种漂移，最终撞壁。若设电场变化尺度 $L$、时间 $\tau$，极化漂移 $\mathbf v_{pol,s}=m_s\dot{\mathbf E}_\perp/(q_sB^2)$；它能在电场形成时造成瞬态极化，不能给无限时间的纯环向静态约束。原题无尺寸／温度数字，不能声称唯一漏失秒数。
\end{solution}

\begin{exercise}\label{ex:hist02-mid-5}
原题5（15分）：轴向磁场的 $\theta$ 箍缩满足 $B_z(r)=B_0(r/a)$（$r<a$）、$B_z=B_0$（$r>a$）。求边缘条件 $p(a)=0$ 的径向压力，再改 $p(a)=p_e$。
\end{exercise}
\begin{solution}
编者独立解答。$\mathbf J=(\nabla\times\mathbf B)/\mu_0$ 仅有 $J_\theta=-(1/\mu_0)\dd B_z/\dd r$，故 $\dd p/\dd r=J_\theta B_z=-(1/2\mu_0)\dd B_z^2/\dd r$，即 $p+B_z^2/(2\mu_0)$ 常量。$r<a$ 时
\begin{equation}p(r)=p_e+\frac{B_0^2}{2\mu_0}\left(1-\frac{r^2}{a^2}\right).
\end{equation}
第一问令 $p_e=0$；第二问在整个内部加 $p_e$，若外区压力为零，则边界需额外表面力或磁场跃变支持非零 $p_e$，不能把内部解和外部真空平衡无条件拼接。
\end{solution}

\begin{exercise}\label{ex:hist02-mid-6}
原题6（16分）：冷电子静止、冷质子沿 $z$ 以 $V_0$ 漂移，说明 $D=1-\omega_{pe}^2/\omega^2-\omega_{pi}^2/(\omega-kV_0)^2$。先令电子响应为零求离子两支及波能符号，再说明恢复电子响应后如何不稳定；附加题估增长率，$m_e/m_i\ll1$。
\end{exercise}
\begin{solution}
编者独立解答。每种冷流 $-\ii(\omega-ku_{0s})\delta v_s=q_s\delta E/m_s$、连续式 $\delta n_s=k n_s\delta v_s/(\omega-ku_{0s})$ 代入 Poisson，得题给 $D$。冻结电子时 $\omega=kV_0\pm\omega_{pi}$；其波能符号由 $\omega\partial_\omega D$ 给出，两支在适当 $k$ 范围有相反符号，允许与电子正能 Langmuir 支耦合增长。取 $kV_0\simeq\omega_{pe}$、$\omega=\omega_{pe}+\delta$，展开 $1-\omega_{pe}^2/\omega^2\simeq2\delta/\omega_{pe}$，有 $2\delta^3/\omega_{pe}\simeq\omega_{pi}^2$。三根中一对共轭，增长根
\begin{equation}\gamma\simeq\frac{\sqrt3}{2}\left(\frac{\omega_{pe}\omega_{pi}^2}{2}\right)^{1/3},\qquad k\simeq\omega_{pe}/V_0.
\end{equation}
只在冷束与近共振、$\gamma\ll\omega_{pe}$ 的排序下成立；有限热宽可稳定窄共振。
\end{solution}
